%%%%%%%%%%%%%%%%
%%% exod 32 %%%
%%%%%%%%%%%%%%%%

\Chapter{32}

\Verse{1}
{
	\hE{וַיַּרְא הָעָם כִּי בֹשֵׁשׁ מֹשֶׁה לָרֶדֶת מִן הָהָר וַיִּקָּהֵל הָעָם עַל אַהֲרֹן וַיֹּאמְרוּ אֵלָיו קוּם עֲשֵׂה לָנוּ אֱלֹהִים אֲשֶׁר יֵלְכוּ לְפָנֵינוּ כִּי זֶה מֹשֶׁה הָאִישׁ אֲשֶׁר הֶעֱלָנוּ מֵאֶרֶץ מִצְרַיִם לֹא יָדַעְנוּ מֶה הָיָה לוֹ׃}
}
{
	\eE{
		32 And the people saw that Moses was delaying to come down\\
		from the mountain, and the people assembled at Aaron and\\
		said to him, “Get up. Make gods for us who will go in\\
		front of us, because this Moses, the man who brought us up\\
		from the land of Egypt: we don’t know what has become of\\
		him!”\\
	}
}
{
%	\footnote{
%		this Moses, the man who brought us up from the land of Egypt. They talk
%		about Moses as if it is necessary to identify who he is. That is, they
%		minimize him and push him into the past tense, to make the case to Aaron
%		that he should act. Footnote 32:1. the man who brought us up from the
%		land of Egypt. Less than forty days earlier they heard the divine voice
%		itself say that God brought them out of Egypt, but, like the wine
%		steward and Joseph’s Pharaoh, they still focus on the human rather than
%		on the deity.
%	}
}

\Verse{2}
{
	\hE{וַיֹּאמֶר אֲלֵהֶם אַהֲרֹן פָּרְקוּ נִזְמֵי הַזָּהָב אֲשֶׁר בְּאזְנֵי נְשֵׁיכֶם בְּנֵיכֶם וּבְנֹתֵיכֶם וְהָבִיאוּ אֵלָי׃}
}
{
	\eE{
		And Aaron said to them, “Take off the gold rings that are\\
		in your wives’, your sons’, and your daughters’ ears, and\\
		bring them to me.”\\
	}
}
{
}

\Verse{3}
{
	\hE{וַיִּתְפָּרְקוּ כּל הָעָם אֶת נִזְמֵי הַזָּהָב אֲשֶׁר בְּאזְנֵיהֶם וַיָּבִיאוּ אֶל אַהֲרֹן׃}
}
{
	\eE{
		And all the people took off the gold rings that were in\\
		their ears and brought them to Aaron.\\
	}
}
{
}

\Verse{4}
{
	\hE{וַיִּקַּח מִיָּדָם וַיָּצַר אֹתוֹ בַּחֶרֶט וַיַּעֲשֵׂהוּ עֵגֶל מַסֵּכָה וַיֹּאמְרוּ אֵלֶּה אֱלֹהֶיךָ יִשְׂרָאֵל אֲשֶׁר הֶעֱלוּךָ מֵאֶרֶץ מִצְרָיִם׃}
}
{
	\eE{
		And he took them from their hand and fashioned it with a\\
		stylus and made it a molten calf. And they said, “These\\
		are your gods, Israel, who brought you up from the land of\\
		Egypt!”\\
	}
}
{
%	\footnote{
%		calf. “Calf” is a bit misleading. It is a young bull. I retained the
%		word “calf” in the translation because the episode is now famous as the
%		“golden calf” story. 32:4–5. These are your gods.... And Aaron saw … A
%		festival to YHWH tomorrow. So many questions: Why do they use the
%		plural, “These are your gods,” when there is only one calf? What did
%		Aaron see? Why does he proclaim a festival to YHWH in the middle of a
%		pagan heresy? I gave the critical explanation of these problems
%		elsewhere (see Who Wrote the Bible?). Here I want to give a pesat
%		explanation of them as they appear in the text. The plural “gods” does
%		not appear to refer to the calf. In the pagan religions of the ancient
%		Near East, people did not worship animals, and they did not worship
%		idols. They worshiped gods who were represented or honored by statues.
%		The gods were sometimes pictured as enthroned on a platform, and the
%		platform was commonly composed of figures of cherubs or of young bulls.
%		Since there is only one young bull here, together with plural “gods,” we
%		must understand the people to be perceiving the bull as a throne
%		platform for gods whom they will worship. Note that the calf itself is
%		not referred to as an idol. What Aaron sees is this behavior on the
%		people’s part. The text does not reveal what Aaron has in mind when he
%		makes the golden calf: does he mean it as a throne platform for such
%		pagan deities, or as a throne platform for YHWH?! All we know is that,
%		after he sees the people associating it with gods and claiming that
%		these gods brought them out of Egypt, he declares that the worship that
%		is to take place is “A festival to YHWH!” (Indeed, the word for “And
%		Aaron saw” is Hebrew . In the original Hebrew text, before the vowels
%		were added, it can also have meant “And Aaron was afraid”! That would
%		make it even clearer why he would redirect the worship back to YHWH.)
%		Aaron is not completely innocent in the matter of the golden calf, as we
%		can see from the defense he offers to Moses afterward (32:24); but it
%		may be that he is to be credited as a person who recognizes that he has
%		made a dangerous mistake and tries to turn it to something good. In
%		recognizing his own error, and in trying to steer a people who have gone
%		wrong onto a different course, Aaron becomes an extraordinarily valuable
%		model—and the story teaches an extraordinarily valuable lesson.
%	}
}

\Verse{5}
{
	\hE{וַיַּרְא אַהֲרֹן וַיִּבֶן מִזְבֵּחַ לְפָנָיו וַיִּקְרָא אַהֲרֹן וַיֹּאמַר חַג לַיהֹוָה מָחָר׃}
}
{
	\eE{
		And Aaron saw, and he built an altar in front of it, and\\
		Aaron called, and he said, “A festival to YHWH tomorrow!”\\
	}
}
{
}

\Verse{6}
{
	\hE{וַיַּשְׁכִּימוּ מִמּחֳרָת וַיַּעֲלוּ עֹלֹת וַיַּגִּשׁוּ שְׁלָמִים וַיֵּשֶׁב הָעָם לֶאֱכֹל וְשָׁתוֹ וַיָּקֻמוּ לְצַחֵק׃}
}
{
	\eE{
		And they got up early the next day, and they made burnt\\
		offerings and brought over peace offerings, and the people\\
		sat down to eat and drink, and they got up to fool around.\\
	}
}
{
%	\footnote{
%		fool around. Some of the occurrences of this word involve sex (Gen 26:8;
%		39:14,17), and some do not (Gen 19:14; 21:9; Judg 16:25). It is not
%		justified, therefore, to assume that sexual play is implied here.
%	}
}

\Verse{7}
{
	\hJ{וַיְדַבֵּר יְהֹוָה אֶל מֹשֶׁה לֶךְ רֵד כִּי שִׁחֵת עַמְּךָ אֲשֶׁר הֶעֱלֵיתָ מֵאֶרֶץ מִצְרָיִם׃}
}
{
	\eJ{
		And YHWH spoke to Moses: “Go. Go down. Because your\\
		people, whom you brought up from the land of Egypt, has\\
		corrupted.\\
	}
}
{
%	\footnote{
%		your people, whom you brought up from the land of Egypt. The issue of
%		who brought the people out, which came up in v. 1, now has an
%		extraordinary twist. Like a parent who says to one’s spouse about a
%		naughty child, “Guess what your child did today!”—so God now says to
%		Moses, “Your people, whom you brought up from the land of Egypt, has
%		corrupted!” But look at Moses’ words in response. He seeks to appeal to
%		God’s compassion and feelings for the people, and so he turns the
%		pronoun around, saying, “Why should your anger flare at your people,
%		whom you brought out from the land of Egypt? … Relent about the bad to
%		your people.” And the text conveys that Moses’ case is successful,
%		because it concludes: “YHWH relented about the bad that He had spoken,
%		to do to His people”!
%	}
}

\Verse{8}
{
	\hJ{סָרוּ מַהֵר מִן הַדֶּרֶךְ אֲשֶׁר צִוִּיתִם עָשׂוּ לָהֶם עֵגֶל מַסֵּכָה וַיִּשְׁתַּחֲווּ לוֹ וַיִּזְבְּחוּ לוֹ וַיֹּאמְרוּ אֵלֶּה אֱלֹהֶיךָ יִשְׂרָאֵל אֲשֶׁר הֶעֱלוּךָ מֵאֶרֶץ מִצְרָיִם׃}
}
{
	\eJ{
		They’ve turned quickly from the way that I commanded them.\\
		They’ve made themselves a molten calf, and they’ve bowed\\
		to it and sacrificed to it and said, ‘These are your gods,\\
		Israel, who brought you up from the land of Egypt!’”\\
	}
}
{
%	\footnote{
%		They’ve turned quickly. How could they have committed the golden calf
%		sin so soon after the revelation—within forty days after actually
%		hearing God speak from the sky? It is when God is closest that humans
%		commit the greatest sin. Similarly in the story of Eden, where God walks
%		among humans, they violate a direct command. Similarly through the
%		wilderness period to come, in which the people see miracles daily
%		(manna, column of cloud and fire), they are repeatedly rebellious. It
%		appears that close proximity to divinity—to an all-seeing, all-knowing
%		divine parent who is always watching, commanding, and judging—is barely
%		tolerable for human beings. Like adolescents who love and need their
%		parents but who long for independence, humankind is pictured as longing
%		to be out of the divine shadow. Note that at the end of the Tanak, in
%		the books of Ezra, Nehemiah, and Esther, there are no miracles, angels,
%		or appearances of God. The words “And God spoke to” never occur. And in
%		these books, where God is the least visible, humans are pictured as
%		relatively good, listening to their leaders, changing their ways when
%		criticized, and observing the commandments.
%	}
}

\Verse{9}
{
	\hJ{וַיֹּאמֶר יְהֹוָה אֶל מֹשֶׁה רָאִיתִי אֶת הָעָם הַזֶּה וְהִנֵּה עַם קְשֵׁה עֹרֶף הוּא׃}
}
{
	\eJ{
		And YHWH said to Moses, “I’ve seen this people; and, here,\\
		it’s a hard-necked people.\\
	}
}
{
}

\Verse{10}
{
	\hJ{וְעַתָּה הַנִּיחָה לִּי וְיִחַר אַפִּי בָהֶם וַאֲכַלֵּם וְאֶעֱשֶׂה אוֹתְךָ לְגוֹי גָּדוֹל׃}
}
{
	\eJ{
		And now, leave off from me, and my anger will flare at\\
		them, and I’ll finish them, and I’ll make you into a big\\
		nation!”\\
	}
}
{
}

\Verse{11}
{
	\hJ{וַיְחַל מֹשֶׁה אֶת פְּנֵי יְהֹוָה אֱלֹהָיו וַיֹּאמֶר לָמָה יְהֹוָה יֶחֱרֶה אַפְּךָ בְּעַמֶּךָ אֲשֶׁר הוֹצֵאתָ מֵאֶרֶץ מִצְרַיִם בְּכֹחַ גָּדוֹל וּבְיָד חֲזָקָה׃}
}
{
	\eJ{
		And Moses conciliated in front of YHWH, his God, and said,\\
		“Why, YHWH, should your anger flare at your people whom\\
		you brought out from the land of Egypt with big power and\\
		with a strong hand?\\
	}
}
{
}

\Verse{12}
{
	\hJ{לָמָּה יֹאמְרוּ מִצְרַיִם לֵאמֹר בְּרָעָה הוֹצִיאָם לַהֲרֹג אֹתָם בֶּהָרִים וּלְכַלֹּתָם מֵעַל פְּנֵי הָאֲדָמָה שׁוּב מֵחֲרוֹן אַפֶּךָ וְהִנָּחֵם עַל הָרָעָה לְעַמֶּךָ׃}
}
{
	\eJ{
		Why should Egypt say, saying, ‘He brought them out for\\
		bad, to kill them in the mountains, and to finish them\\
		from on the face of the earth’? Turn back from your\\
		flaring anger, and relent about the bad to your people.\\
	}
}
{
}

\Verse{13}
{
	\hJ{זְכֹר לְאַבְרָהָם לְיִצְחָק וּלְיִשְׂרָאֵל עֲבָדֶיךָ אֲשֶׁר נִשְׁבַּעְתָּ לָהֶם בָּךְ וַתְּדַבֵּר אֲלֵהֶם אַרְבֶּה אֶת זַרְעֲכֶם כְּכוֹכְבֵי הַשָּׁמָיִם וְכל הָאָרֶץ הַזֹּאת אֲשֶׁר אָמַרְתִּי אֶתֵּן לְזַרְעֲכֶם וְנָחֲלוּ לְעֹלָם׃}
}
{
	\eJ{
		Remember Abraham, Isaac, and Israel, your servants, that\\
		you swore to them by yourself, and you spoke to them:\\
		‘I’ll multiply your seed like the stars of the skies, and\\
		I’ll give to your seed all this land that I’ve said, and\\
		they’ll possess it forever.’”\\
	}
}
{
}

\Verse{14}
{
	\hJ{וַיִּנָּחֶם יְהֹוָה עַל הָרָעָה אֲשֶׁר דִּבֶּר לַעֲשׂוֹת לְעַמּוֹ׃}
}
{
	\eJ{
		And YHWH relented about the bad that He had spoken, to do\\
		to His people.\\
	}
}
{
}

\Verse{15}
{
	\hE{וַיִּפֶן וַיֵּרֶד מֹשֶׁה מִן הָהָר וּשְׁנֵי לֻחֹת הָעֵדֻת בְּיָדוֹ לֻחֹת כְּתֻבִים מִשְּׁנֵי עֶבְרֵיהֶם מִזֶּה וּמִזֶּה הֵם כְּתֻבִים׃}
}
{
	\eE{
		And Moses turned and went down from the mountain. And the\\
		two tablets of witness were in his hand, tablets written\\
		from their two sides: from this side and from this side\\
		they were written.\\
	}
}
{
%	\footnote{
%		tablets of witness. It is a crucial concept here and in Deuteronomy (and
%		Kings) that there is written witness/testimony confirming what God has
%		said, so subsequent generations can never say, “We didn’t know.” Even
%		though the people have actually heard it from the voice of God, the
%		Sinai covenant is a legal procedure and requires a formal, written
%		document. The Noahic and Abrahamic covenants are sealed with a divine
%		oath; and, since God’s violating His own promise is unthinkable, no
%		document is required. But the Sinai covenant is sealed by the people’s
%		promise that “We’ll do it, and we’ll listen”; and human commitments
%		require written witness. That is why there are tablets, and that is why
%		they are called the tablets of witness. Footnote 32:15. written from
%		their two sides. Does this mean: written on both tablets? or written on
%		both the front and the back of each tablet? The former seems too obvious
%		to require this mention: of course it is written on both tablets;
%		otherwise, why have two tablets? So the point here is precisely that
%		they are written on both sides of each tablet. For pictorial purposes
%		(in art and film), it is preferable to put the commandments all on the
%		same side of each tablet, but the text suggests otherwise. And we know
%		of numerous ancient royal inscriptions and statues that were written on
%		both sides, front and back.
%	}
}

\Verse{16}
{
	\hE{וְהַלֻּחֹת מַעֲשֵׂה אֱלֹהִים הֵמָּה וְהַמִּכְתָּב מִכְתַּב אֱלֹהִים הוּא חָרוּת עַל הַלֻּחֹת׃}
}
{
	\eE{
		And the tablets: they were God’s doing. And the writing:\\
		it was God’s writing, inscribed on the tablets.\\
	}
}
{
}

\Verse{17}
{
	\hE{וַיִּשְׁמַע יְהוֹשֻׁעַ אֶת קוֹל הָעָם בְּרֵעֹה וַיֹּאמֶר אֶל מֹשֶׁה קוֹל מִלְחָמָה בַּמַּחֲנֶה׃}
}
{
	\eE{
		And Joshua heard the sound of the people in its shouting,\\
		and he said to Moses, “A sound of war is in the camp.”\\
	}
}
{
}

\Verse{18}
{
	\hE{וַיֹּאמֶר אֵין קוֹל עֲנוֹת גְּבוּרָה וְאֵין קוֹל עֲנוֹת חֲלוּשָׁה קוֹל עַנּוֹת אָנֹכִי שֹׁמֵעַ׃}
}
{
	\eE{
		And he said, “It’s not a sound of singing of victory, and\\
		it’s not a sound of singing of defeat. It’s just the sound\\
		of singing I hear!”\\
	}
}
{
}

\Verse{19}
{
	\hE{וַיְהִי כַּאֲשֶׁר קָרַב אֶל הַמַּחֲנֶה וַיַּרְא אֶת הָעֵגֶל וּמְחֹלֹת וַיִּחַר אַף מֹשֶׁה וַיַּשְׁלֵךְ מִיָּדָו אֶת הַלֻּחֹת וַיְשַׁבֵּר אֹתָם תַּחַת הָהָר׃}
}
{
	\eE{
		And it was when he came close to the camp, and he saw the\\
		calf and dancing: and Moses’ anger flared, and he threw\\
		the tablets from his hands and shattered them below the\\
		mountain.\\
	}
}
{
%	\footnote{
%		Moses’ anger flared. God had said, “My anger will flare at them,” but
%		Moses dissuaded God, saying, “Why should your anger flare at your
%		people?” But when Moses himself returns to the people and sees the bull,
%		the text reports that “his anger flared,” the very words for what he
%		persuaded God not to do! He smashes the tablets that he was bringing
%		back to the people, and he calls for a bloody purge of thousands among
%		the people. Quintessentially among prophets, Moses comes to God as a
%		human, pleading the human case; and he comes to the people as a
%		confidant of God, feeling the divine disappointment and anger with
%		humans (albeit apparently still lacking the divine restraint and
%		capacity to relent). It is the curse of the prophet: to God he
%		represents humans, and to humans he represents God.
%	}
}

\Verse{20}
{
	\hE{וַיִּקַּח אֶת הָעֵגֶל אֲשֶׁר עָשׂוּ וַיִּשְׂרֹף בָּאֵשׁ וַיִּטְחַן עַד אֲשֶׁר דָּק וַיִּזֶר עַל פְּנֵי הַמַּיִם וַיַּשְׁקְ אֶת בְּנֵי יִשְׂרָאֵל׃}
}
{
	\eE{
		And he took the calf that they had made, and he burned it\\
		in fire and ground it until it was thin, and he scattered\\
		it on the face of the water, and he made the children of\\
		Israel drink!\\
	}
}
{
%	\footnote{
%		thin. In Deut 9:21, Moses elaborates that this means “thin as dust.”
%		Many English translations say “ground it to powder.” Footnote 32:20. he
%		scattered it on the face of the water. What water? It is natural to find
%		streams at the bases of mountains, but in fact the last water mentioned
%		is the water that came out of the crag of Horeb/Sinai. That would add
%		irony: that the water that ended one rebellion now carries off the
%		product of the next rebellion. And, as the people now see the dust of
%		the golden calf on the water, they may be reminded that this water came
%		to them from God and through Moses, whom they have been quick to forget.
%		Footnote 32:20. he made the children of Israel drink. Classical and
%		recent commentators have compared this to the law of a woman suspected
%		of adultery (sh), who must drink a mixture of dust and water (Num
%		5:11–31). Indeed, both cases use the unusual word pr‘ (see comment on
%		32:25 below). And the issue in both is betrayal. The prophets, most
%		notably Hosea, compare the people’s turning to other gods with adultery.
%		Still, I urge great caution in interpreting the connection. The people
%		in Exodus are guilty! The sh case is a process designed to determine
%		whether someone is guilty or innocent. And the dust there comes from the
%		Tabernacle, which is something positive, while the dust of the calf here
%		is the opposite. The similarities between the two matters may reflect
%		only very general common concerns.
%	}
}

\Verse{21}
{
	\hE{וַיֹּאמֶר מֹשֶׁה אֶל אַהֲרֹן מֶה עָשָׂה לְךָ הָעָם הַזֶּה כִּי הֵבֵאתָ עָלָיו חֲטָאָה גְדֹלָה׃}
}
{
	\eE{
		And Moses said to Aaron, “What did this people do to you,\\
		that you’ve brought a big sin on it?!”\\
	}
}
{
}

\Verse{22}
{
	\hE{וַיֹּאמֶר אַהֲרֹן אַל יִחַר אַף אֲדֹנִי אַתָּה יָדַעְתָּ אֶת הָעָם כִּי בְרָע הוּא׃}
}
{
	\eE{
		And Aaron said, “Let my lord’s anger not flare. You know\\
		the people, that it’s in a bad state.\\
	}
}
{
%	\footnote{
%		in a bad state. In the sense of “in for trouble,” “up to no good,” or
%		“in a sorry state.”
%	}
}

\Verse{23}
{
	\hE{וַיֹּאמְרוּ לִי עֲשֵׂה לָנוּ אֱלֹהִים אֲשֶׁר יֵלְכוּ לְפָנֵינוּ כִּי זֶה מֹשֶׁה הָאִישׁ אֲשֶׁר הֶעֱלָנוּ מֵאֶרֶץ מִצְרַיִם לֹא יָדַעְנוּ מֶה הָיָה לוֹ׃}
}
{
	\eE{
		And they said to me, ‘Make gods for us who will go in\\
		front of us, because this Moses, the man who brought us up\\
		from the land of Egypt: we don’t know what has become of\\
		him!’\\
	}
}
{
}

\Verse{24}
{
	\hE{וָאֹמַר לָהֶם לְמִי זָהָב הִתְפָּרָקוּ וַיִּתְּנוּ לִי וָאַשְׁלִכֵהוּ בָאֵשׁ וַיֵּצֵא הָעֵגֶל הַזֶּה׃}
}
{
	\eE{
		And I said to them, ‘Whoever has gold: take it off.’ And\\
		they gave it to me, and I threw it into the fire, and this\\
		calf came out!”\\
	}
}
{
%	\footnote{
%		and this calf came out. A strong lie uses as much of the truth as
%		possible. Aaron quotes the people exactly. Then he quotes his own words
%		with slight changes. And then he transforms the production of the calf
%		from his own fashioning to “and this calf came out!” He also addresses
%		Moses as “my lord,” and he emphasizes the people’s bad state. Thus, in a
%		very brief response he does a great deal to minimize any culpability of
%		his own. Why is there no consequence for Aaron? Some say that the
%		consequence is the deaths of his sons Nadab and Abihu (Leviticus 10).
%		Alternatively, the answer may be that just prior to the making of the
%		calf God has designated Aaron as the people’s high priest. In that
%		position he cannot suffer any direct punishment from God, which would
%		both disqualify him and demean the office of high priest. And so what is
%		left is for Aaron to make amends through personal atonement. (The same
%		would apply to the case in which he and Miriam criticize Moses about his
%		wife. Miriam becomes leprous, but Aaron does not suffer any consequence.
%		See the comment on Num 12:11.)
%	}
}

\Verse{25}
{
	\hE{וַיַּרְא מֹשֶׁה אֶת הָעָם כִּי פָרֻעַ הוּא כִּי פְרָעֹה אַהֲרֹן לְשִׁמְצָה בְּקָמֵיהֶם׃}
}
{
	\eE{
		And Moses saw the people, that it was turned loose,\\
		because Aaron had turned it loose—to denigration among\\
		their adversaries!\\
	}
}
{
%	\footnote{
%		turned it loose. The word in Hebrew consonants is pr‘h, a pun on the
%		word “Pharaoh.” Aaron has “Pharaohed” the people; he has done something
%		to them that Pharaohs had done: made them ignoble in the eyes of those
%		who oppose them. Or one could say that he has brought them back to the
%		condition in which they were before the Sinai revelation: in disarray,
%		without the law. Footnote 32:25. denigration. Scornful whispering.
%	}
}

\Verse{26}
{
	\hE{וַיַּעֲמֹד מֹשֶׁה בְּשַׁעַר הַמַּחֲנֶה וַיֹּאמֶר מִי לַיהֹוָה אֵלָי וַיֵּאָסְפוּ אֵלָיו כּל בְּנֵי לֵוִי׃}
}
{
	\eE{
		And Moses stood in the gate of the camp and said, “Whoever\\
		is for YHWH: to me!” And all the children of Levi were\\
		gathered to him.\\
	}
}
{
}

\Verse{27}
{
	\hE{וַיֹּאמֶר לָהֶם כֹּה אָמַר יְהֹוָה אֱלֹהֵי יִשְׂרָאֵל שִׂימוּ אִישׁ חַרְבּוֹ עַל יְרֵכוֹ עִבְרוּ וָשׁוּבוּ מִשַּׁעַר לָשַׁעַר בַּמַּחֲנֶה וְהִרְגוּ אִישׁ אֶת אָחִיו וְאִישׁ אֶת רֵעֵהוּ וְאִישׁ אֶת קְרֹבוֹ׃}
}
{
	\eE{
		And he said to them, “YHWH, God of Israel, said this:\\
		‘Set, each man, his sword on his thigh; cross over and\\
		come back from gate to gate in the camp; and kill, each\\
		man, his brother and, each man, his neighbor and, each\\
		man, his relative.’”\\
	}
}
{
}

\Verse{28}
{
	\hE{וַיַּעֲשׂוּ בְנֵי לֵוִי כִּדְבַר מֹשֶׁה וַיִּפֹּל מִן הָעָם בַּיּוֹם הַהוּא כִּשְׁלֹשֶׁת אַלְפֵי אִישׁ׃}
}
{
	\eE{
		And the children of Levi did according to Moses’ word, and\\
		about three thousand men fell from the people in that day.\\
	}
}
{
}

\Verse{29}
{
	\hE{וַיֹּאמֶר מֹשֶׁה מִלְאוּ יֶדְכֶם הַיּוֹם לַיהֹוָה כִּי אִישׁ בִּבְנוֹ וּבְאָחִיו וְלָתֵת עֲלֵיכֶם הַיּוֹם בְּרָכָה׃}
}
{
	\eE{
		And Moses said, “Fill your hand to YHWH today—because each\\
		man was at his son and at his brother—and to put a\\
		blessing on you today.”\\
	}
}
{
%	\footnote{
%		Fill your hand. This phrase refers to volunteering for divine service.
%		See the comment on 28:41. Footnote 32:29. to YHWH today—because each man
%		was at his son and at his brother—and to put a blessing on you today.
%		The grammar of this verse is difficult in the Hebrew and more so in
%		translation. It appears to mean that, because of the Levites’ violent
%		zeal, they are now chosen to consecrate themselves as Israel’s
%		priesthood and are to receive divine blessing. It is ironic that Levi is
%		criticized in Jacob’s deathbed blessing for his violent slaughter in
%		Shechem (Gen 34:25; 49:5–7), but Levi’s descendants are rewarded for a
%		violent slaughter here. The lesson may be that an act is not judged, in
%		the Torah, apart from its circumstances. What the Levites do in response
%		to a heresy of destiny-changing proportions is not comparable to Levi’s
%		(and Simeon’s) massacre of a city in response to the offense committed
%		by their ruler’s son.
%	}
}

\Verse{30}
{
	\hJ{וַיְהִי מִמּחֳרָת וַיֹּאמֶר מֹשֶׁה אֶל הָעָם אַתֶּם חֲטָאתֶם חֲטָאָה גְדֹלָה וְעַתָּה אֶעֱלֶה אֶל יְהֹוָה אוּלַי אֲכַפְּרָה בְּעַד חַטַּאתְכֶם׃}
}
{
	\eJ{
		And it was on the next day, and Moses said to the people,\\
		“You’ve committed a big sin. And now I’ll go up to YHWH.\\
		Perhaps I may make atonement for your sin.”\\
	}
}
{
}

\Verse{31}
{
	\hJ{וַיָּשׁב מֹשֶׁה אֶל יְהֹוָה וַיֹּאמַר אָנָּא חָטָא הָעָם הַזֶּה חֲטָאָה גְדֹלָה וַיַּעֲשׂוּ לָהֶם אֱלֹהֵי זָהָב׃}
}
{
	\eJ{
		And Moses went back to YHWH and said, “Please, this people\\
		has committed a big sin and made gods of gold for\\
		themselves.\\
	}
}
{
%	\footnote{
%		Please. This word, Hebrew ’nnâ, occurs only twice in the Torah, here and
%		in Gen 50:17; and in both instances it is used in the context of asking
%		someone to “bear a sin.” In both cases it is a tremendous sin, and in
%		both cases there are no grounds for defense. It is thus a particle that
%		expresses emotional entreaty. Interestingly, the case in Genesis 50 is
%		the plea by Joseph’s brothers to forgive them for what they did to him,
%		and Joseph does in fact forgive them. Again the story of Moses and the
%		Israelites contains an allusion to a story from the era of the
%		patriarchs, hinting to the reader that there is reason for hope; for if
%		a great man can forgive, certainly God can as well. Footnote 32:31. gods
%		of gold. Or “a god of gold.” It is uncertain whether Moses refers here
%		to the calf statue as a god or if he refers to gods for whom the calf is
%		a throne dais. Footnote 32:31. made gods of gold for themselves. Moses
%		does not describe what they have done in the language of the Ten
%		Commandments (“You shall not make a statue or any form … ”); he rather
%		uses the wording of one of the commands that come later: “You shall not
%		make gods of gold for yourselves” (Exod 20:20). Breaking one of the Ten
%		Commandments is far more serious than a violation of the other
%		commandments of the Torah, because the Ten Commandments form the text of
%		the Sinai covenant, and to violate them is to breach the covenant
%		itself. There are cases in the Tanak in which individuals or the nation
%		are faced with a choice of breaking one of the Ten Commandments or else
%		one of the other laws. They choose to break the other laws, even though
%		there are terrible consequences, rather than break one of the Ten.
%	}
}

\Verse{32}
{
	\hJ{וְעַתָּה אִם תִּשָּׂא חַטָּאתָם וְאִם אַיִן מְחֵנִי נָא מִסִּפְרְךָ אֲשֶׁר כָּתָבְתָּ׃}
}
{
	\eJ{
		And now, if you will bear their sin—and if not, wipe me\\
		out from your scroll that you’ve written.”\\
	}
}
{
%	\footnote{
%		if you will bear their sin—and if not … Moses seems to break his
%		sentence in the middle and jump to the other side of the coin: “If you
%		will bear their sin”—then what? Perhaps there is nothing to say. He is
%		not in a position to offer God a reward if He does bear their sin. He
%		can only turn to the alternative: “If you won’t bear their sin, then
%		wipe me out from your scroll …” Footnote 32:32. your scroll that you’ve
%		written. The reference is either to the scroll of the Covenant Code
%		(Exod 24:7) or to the scroll concerning Amalek (17:14). In the account
%		of the Amalek scroll, the same word that is used here for “to wipe out”
%		occurs twice. Moses may be implying, “Please do this, or else erase me
%		from history the way you said you would erase Amalek.”
%	}
}

\Verse{33}
{
	\hJ{וַיֹּאמֶר יְהֹוָה אֶל מֹשֶׁה מִי אֲשֶׁר חָטָא לִי אֶמְחֶנּוּ מִסִּפְרִי׃}
}
{
	\eJ{
		And YHWH said to Moses, “The one who has sinned against\\
		me, I’ll wipe him out from my scroll. 32:33–34. The one\\
		who has sinned against me, I’ll wipe him out.... And now,\\
		go. Moses’ gesture is both selfless and audacious at the\\
		same time. He takes a great risk to himself for the sake\\
		of winning forgiveness for his people, but at the same\\
		time he is speaking more boldly to God than he has dared\\
		to do before—and he is ascribing much importance to\\
		himself in assuming that “wipe me out of your scroll”\\
		carries enough weight to affect a divine decision. But\\
		God’s answer seems to brush aside Moses’ gesture, as if to\\
		say, “Don’t tell me whom to write in my scroll. I’ll do\\
		whatever I’ll do. Now, go.”\\
	}
}
{
}

\Verse{34}
{
	\hJ{וְעַתָּה לֵךְ נְחֵה אֶת הָעָם אֶל אֲשֶׁר דִּבַּרְתִּי לָךְ הִנֵּה מַלְאָכִי יֵלֵךְ לְפָנֶיךָ וּבְיוֹם פּקְדִי וּפָקַדְתִּי עֲלֵהֶם חַטָּאתָם׃}
}
{
	\eJ{
		And now, go. Lead the people to where I spoke to you.\\
		Here, my angel will go ahead of you. And, in the day that\\
		I take account, I’ll account their sin on them.”\\
	}
}
{
}

\Verse{35}
{
	\hJ{וַיִּגֹּף יְהֹוָה אֶת הָעָם עַל אֲשֶׁר עָשׂוּ אֶת הָעֵגֶל אֲשֶׁר עָשָׂה אַהֲרֹן׃}
}
{
	\eJ{
		And YHWH struck the people because they had made the calf,\\
		which Aaron had made.\\
	}
}
{
%	\footnote{
%		they had made the calf, which Aaron had made. Some have taken this
%		wording to be a scribal error because it says both that the people made
%		it and that Aaron made it. (And some of the versions say: “they had
%		served the calf which Aaron had made.”) But both did make it, and both
%		are culpable. The community produced the calf: initiating the event,
%		asking for gods to be made, contributing the gold, responding to the
%		finished product. And Aaron was responsible in two ways: as the leader
%		of the community and as the one who fashioned the calf. A group and an
%		individual who acts as the group’s agent share in the responsibility for
%		the things they do.
%	}
}
